% संपूर्ण मराठी ग्रंथातील प्रकरण 16: प्रथम-क्रम तर्कशास्त्राची चिन्हार्थमीमांसा
% हा संपादनयोग्य स्रोतखंड आहे. संकलनासाठी संपूर्ण स्रोतसंग्रहातील
% अक्षररूपे, पूर्वभाग, चिन्हव्याख्या आणि आकृती-स्रोत आवश्यक आहेत.
% मूळ: Open Logic Project; मजकूर आणि रूपांतर CC BY 4.0.
% यंत्रानुवाद, दुरुस्ती आणि तपासणी: OpenAI Codex —
% GPT-5.6 Sol आणि GPT-6 Sol, दोन्ही Ultra effort.
% दुरुस्ती-पुनरावलोकन: GPT-6.1 Sol, Ultra effort.
\chapter{प्रथम-क्रम तर्कशास्त्राची चिन्हार्थमीमांसा}\label{fol:sem::chap}









\section{परिचय}\label{fol:syn:its:sec}

व्यंजकांना अर्थ देणे हे चिन्हार्थमीमांसेचे क्षेत्र आहे. चिन्हार्थमीमांसेतील
मध्यवर्ती संकल्पना म्हणजे संरचना मधील पूर्ति. संरचना
भाषेच्या मूलभूत घटकांना अर्थ देते: प्रांत म्हणजे वस्तूंचा अरिक्त
संच. संख्यापक या प्रांतातील घटकांवर फिरतात असा त्यांचा अर्थ लावला जातो,
स्थिरांक ना प्रांतातील घटक नेमून दिले जातात, फलने ना
प्रांत पासून त्याच प्रांताकडे जाणारी फलने नेमून दिली जातात आणि
विधेये ना प्रांत वरील संबंध नेमून दिले जातात. प्रांत
आणि मूलभूत शब्दसंग्रहातील चिन्हांना केलेल्या नेमणुका मिळून
संरचना बनते.
चले ही सूत्रे मध्ये येऊ शकतात; त्यांना चिन्हार्थ देण्यासाठी
आपल्याला प्रांत मधील घटक ही नेमून द्यावी लागतात---हेच
चर-मूल्यांकन होय. शेवटी, पूर्ति-संबंध या सर्व गोष्टी एकत्र आणतो.
सूत्राची संरचना~$\Struct{M}$ मध्ये चल च्या
मूल्यांकन~$s$ च्या सापेक्ष पूर्ति होऊ शकते; हे $\Sat{M}{\varphi}[s]$ असे
लिहितात. हा संबंधदेखील~$\varphi$ च्या संरचनेवरील विगमनाने परिभाषित केला जातो.
उदाहरणार्थ, $(\varphi \land \psi)$ ची पूर्ति परिभाषित करताना तार्किक संयोजकांचे
सत्यता-कोष्टक वापरून ती $\varphi$ आणि~$\psi$ यांची पूर्ति होते (किंवा होत नाही)
यांच्या आधारे सांगितली जाते. त्यानंतर असे निष्पन्न होते की सूत्र~$\varphi$
हे वाक्य, म्हणजे त्यात मुक्त चरे नसतील, तर चल चे
मूल्यांकन असंबद्ध ठरते; त्यामुळे संरचनांमध्ये वाक्यांची
सरळ पूर्ति होते (किंवा होत नाही) असे आपण म्हणू शकतो.

वाक्ये साठीच्या पूर्ति-संबंध~$\Sat{M}{\varphi}$ च्या आधारे आपण वैधता,
तार्किक निष्पन्नता आणि पूर्ततायोग्यता या मूलभूत चिन्हार्थविषयक संकल्पना
परिभाषित करू शकतो. वाक्याची प्रत्येक संरचनांमध्ये पूर्ति
होत असेल, तर ते वैध, $\Entails \varphi$, असते. वाक्यांच्या एखाद्या
संचातून ते तार्किकरीत्या निष्पन्न होते, $\Gamma \Entails \varphi$, जर~$\Gamma$ मधील
सर्व वाक्यांची पूर्ति करणारे प्रत्येक संरचना हे~$\varphi$ चीही
पूर्ति करत असेल. आणि एखाद्या वाक्यांच्या संचातील सर्व
वाक्यांची एकाच वेळी पूर्ति करणारे एखादे संरचना असेल, तर तो
संच पूर्ततायोग्य असतो. सूत्रे विगमनाधारित रीतीने परिभाषित केली आहेत
आणि पूर्ति हीही सूत्रांच्या संरचनेवरील विगमनाने परिभाषित केली जाते;
म्हणून आपल्या चिन्हार्थमीमांसेचे गुणधर्म सिद्ध करण्यासाठी आणि परिभाषित
चिन्हार्थविषयक संकल्पनांचा परस्परसंबंध लावण्यासाठी विगमन वापरता येते.












\section{प्रथम-क्रम भाषांसाठीच्या संरचना}\label{fol:syn:str:sec}

\begin{explain}
प्रथम-क्रम भाषांना स्वतःचा कोणताही \emph{अर्थ लावलेला नसतो:}
स्थिरांक, फलने आणि विधेये यांना कोणताही विशिष्ट
अर्थ जोडलेला नसतो. एखादी \emph{संरचना}
विनिर्दिष्ट करून अर्थ दिले जातात. ती \emph{प्रांत}, म्हणजे
स्थिरांक ज्या वस्तू निर्देशित करतात, फलने ज्यांवर क्रिया
करतात आणि संख्यापक ज्या वस्तूंवर फिरतात त्या वस्तू विनिर्दिष्ट करते.
याशिवाय, कोणते स्थिरांक कोणत्या वस्तू निर्देशित करतात,
फलन वस्तूंचे वस्तूंमध्ये प्रतिचित्रण कसे करते आणि विधेये
कोणत्या वस्तूंना लागू होतात हे ती विनिर्दिष्ट करते. संरचना या
तर्कशास्त्रातील \emph{चिन्हार्थविषयक} संकल्पनांचा---उदाहरणार्थ, तार्किक
निष्पन्नता, वैधता आणि पूर्ततायोग्यता या संकल्पनांचा---आधार आहेत. साहित्यात त्यांना
निरनिराळ्या ठिकाणी ``रचना,'' ``अर्थनिर्धारण'' किंवा ``प्रतिमान'' असे
म्हटले जाते.
\end{explain}

\begin{defn}[संरचना]
प्रथम-क्रम तर्कशास्त्राची भाषा~$\Lang{L}$ हिच्यासाठीची
\emph{संरचना}~$\Struct M$ पुढील घटकांनी बनते:
\begin{enumerate}
\item \emph{प्रांत:} अरिक्त संच~$\Domain M$.
\item \emph{स्थिरांकाचे अर्थनिर्धारण:} $\Lang{L}$ मधील प्रत्येक
  स्थिरांक~$c$ साठी घटक~$\Assign{c}{M} \in \Domain M$.
\item \emph{विधेयांचे अर्थनिर्धारण:} $\Lang{L}$ मधील प्रत्येक
  $n$-स्थानी विधेय~$R$ साठी ($\eq$ सोडून) $n$-स्थानी संबंध
  $\Assign{R}{M} \subseteq \Domain{M}^n$.
\item \emph{फलनांचे अर्थनिर्धारण:} $\Lang{L}$ मधील प्रत्येक
  $n$-स्थानी फलन~$f$ साठी $n$-स्थानी फलन
  $\Assign{f}{M} \colon \Domain{M}^n \to \Domain{M}$.
\end{enumerate}
\end{defn}

\begin{ex}
अंकगणिताच्या भाषेसाठीच्या संरचना~$\Struct M$ मध्ये एक संच,
स्थिरांक~$\Obj 0$ चे अर्थनिर्धारण म्हणून $\Domain M$ मधील
घटक~$\Assign{\Obj 0}{M}$, एकस्थानी फलन
$\Assign{\Obj \prime}{M} \colon \Domain{M} \to \Domain M$,
$\Domain M^2 \to \Domain M$ अशी दोन द्विस्थानी फलने
$\Assign{\Obj +}{M}$ आणि $\Assign{\Obj \times}{M}$, तसेच द्विस्थानी
संबंध~$\Assign{\Obj <}{M} \subseteq \Domain{M}^2$ असतात.

अशा रचनेचे एक स्वाभाविक उदाहरण पुढीलप्रमाणे आहे:
\begin{enumerate}
\item $\Domain N = \Nat$
\item $\Assign{\Obj 0}{N} = 0$
\item प्रत्येक $n \in \Nat$ साठी $\Assign{\Obj \prime}{N}(n) = n + 1$
\item प्रत्येक $n, m \in \Nat$ साठी $\Assign{\Obj +}{N}(n, m) = n + m$
\item प्रत्येक $n, m \in \Nat$ साठी $\Assign{\Obj \times}{N}(n, m) = n\cdot m$
\item $\Assign{\Obj <}{N} = \Setabs{\tuple{n, m}}{n \in \Nat, m \in
  \Nat, n < m}$
\end{enumerate}
अशा रीतीने परिभाषित केलेल्या $\Lang L_A$ साठीच्या रचनेला~$\Struct N$
\emph{अंकगणिताचे मानक प्रतिमान} म्हणतात, कारण ती~$\Lang L_A$ मधील
अतार्किक चिन्हांचा अर्थ आपण अपेक्षा करू त्याच रीतीने लावते.

तरीही, $\Lang L_A$ साठी इतर अनेक संरचना शक्य आहेत. उदाहरणार्थ,
प्रांत म्हणून~$\Nat$ ऐवजी पूर्णांकांचा संच~$\Int$ घेऊन $\Obj 0$,
$\Obj \prime$, $\Obj +$, $\Obj \times$, $\Obj <$ यांची अर्थनिर्धारणे
त्यानुसार परिभाषित करता येतील. पण संख्यांशी दुरान्वयानेही संबंध नसलेल्या
$\Lang L_A$ साठीच्या रचनाही आपण परिभाषित करू शकतो.
\end{ex}

\begin{ex}
संचसिद्धान्ताची भाषा~$\Lang L_Z$ हिच्यासाठीची रचना~$\Struct M$ हिला
फक्त एक संच आणि एकच द्विस्थानी संबंध आवश्यक असतात. म्हणून तांत्रिकदृष्ट्या,
उदाहरणार्थ, माणसांचा संच आणि ``$x$ हा~$y$ पेक्षा वयाने मोठा आहे'' हा संबंध
$\Lang L_Z$ साठी संरचना म्हणून वापरता येईल; त्याचप्रमाणे~$\Nat$
आणि $n, m \in \Nat$ साठीचा $n \ge m$ हाही वापरता येईल.


$\Lang L_Z$ साठीची एक विशेष रोचक संरचना, जिच्या प्रांतातील
घटक प्रत्यक्षात संच आहेत आणि जिच्यात $\Obj \in$ चे अर्थनिर्धारण
प्रत्यक्षात ``$x$ हा~$y$ चा घटक आहे'' हा संबंध आहे, ती म्हणजे
\emph{आनुवंशिकदृष्ट्या सांत संचांची} संरचना~$\Struct{{HF}}$:
\begin{enumerate}
\item $\Domain{{{HF}}} = \emptyset \cup \Pow{\emptyset} \cup
  \Pow{\Pow{\emptyset}} \cup \Pow{\Pow{\Pow{\emptyset}}} \cup \dots$;
\item $\Assign{\Obj \in}{{{HF}}} = \Setabs{\tuple{x, y}}{x, y \in
  \Domain{{{HF}}}, x \in y}$.
\end{enumerate}
\end{ex}

\begin{digress}
संरचना म्हणून काय गणले जाईल याविषयी आपण केलेल्या अटींचा आपल्या
तर्कशास्त्रावर परिणाम होतो. उदाहरणार्थ, प्रांत अरिक्त ठेवण्याची निवड
संख्यकीकृत वाक्यांच्या पूर्तिची (किंवा सत्यतेची) नेहमीची मांडणी गृहीत
धरल्यास $\lexists[x][(\varphi(x) \lor \lnot \varphi(x))]$ वैध---म्हणजे तार्किक
सत्य---असल्याची खात्री देते. तसेच, सर्व स्थिरांक नी प्रांतातील
वस्तूच निर्देशित केली पाहिजे ही अट अस्तित्ववाची सामान्यीकरण हा
निर्दोष अनुमान-रूपबंध असल्याची खात्री देते: $\varphi(a)$, म्हणून
$\lexists[x][\varphi(x)]$. नावांना प्रांताबाहेरील वस्तू निर्देशित करण्याची
किंवा काहीच निर्देशित न करण्याची मुभा दिली, तर आपण \emph{मुक्त
तर्कशास्त्राच्या} दिशेने जाऊ; त्यात अस्तित्ववाची सामान्यीकरणासाठी एक
अतिरिक्त आधारवाक्यालागते: $\varphi(a)$ आणि $\lexists[x][\eq[x][a]]$, म्हणून
$\lexists[x][\varphi(x)]$.
\end{digress}












\section{प्रथम-क्रम भाषांसाठी बंद पदांनी आच्छादित संरचना}\label{fol:syn:cov:sec}

\begin{explain}
एखाद्या पदात चले नसतील, तर ते \emph{बंद} असते, हे आठवा.
\end{explain}

\begin{defn}[बंद पदांचे मूल्य]
$t$ हे भाषा~$\Lang L$ हिचे बंद पद आणि $\Struct M$ ही~$\Lang L$ साठीची
संरचना असेल, तर \emph{मूल्य}~$\Value{t}{M}$ पुढीलप्रमाणे
परिभाषित केले जाते:
\begin{enumerate}
\item $t$ हे फक्त स्थिरांक~$c$ असेल, तर
  $\Value{c}{M} = \Assign{c}{M}$.
\item $t$ हे $\Atom{f}{t_1, \ldots, t_n}$ या रूपात असेल, तर
  \[  \Value{t}{M} = \Assign{f}{M}(\Value{t_1}{M}, \ldots,
  \Value{t_n}{M}).
  \]
\end{enumerate}
\end{defn}

\begin{defn}[आच्छादित संरचना]
प्रांतातील प्रत्येक घटक हा एखाद्या बंद पदाचे मूल्य असेल, तर ती
संरचना \emph{आच्छादित} असते.
\end{defn}

\begin{ex}
स्थिरांक $\Obj{zero}$, $\Obj{one}$, $\Obj{two}$, \dots, द्विस्थानी
विधेय~$<$ आणि द्विस्थानी फलने $+$ व~$\times$ असलेली
भाषा~$\Lang L$ घेऊ. मग~$\Lang L$ साठीची संरचना~$\Struct M$ अशी
आहे: तिचा प्रांत $\Domain M = \{0, 1, 2, \ldots \}$ असून
$\Assign{\Obj{zero}}{M} = 0$, $\Assign{\Obj{one}}{M} = 1$,
$\Assign{\Obj{two}}{M} = 2$ आणि असेच पुढे ही तिची अर्थनिर्धारणे आहेत.
द्विस्थानी संबंधचिन्ह~$<$ साठी $\Assign{<}{M}$ हा अशा सर्व जोड्यांचा
संच आहे की $\tuple{c_1, c_2} \in \Domain{M}^2$ आणि $c_1$ हा~$c_2$ पेक्षा
लहान आहे: उदाहरणार्थ, $\tuple{1, 3} \in \Assign{<}{M}$, पण
$\tuple{2, 2} \notin \Assign{<}{M}$. द्विस्थानी फलन~$+$ साठी
$\Assign{+}{M}$ नेहमीच्या रीतीने परिभाषित करू---उदाहरणार्थ,
$\Assign{+}{M}(2,3)$ हे~$5$ कडे पाठवते; आणि द्विस्थानी
फलन~$\times$ साठीही त्याचप्रमाणे. म्हणून $\Obj{four}$ चे
मूल्य फक्त~$4$ आहे आणि $\times(\Obj{two},
+(\Obj{three},\Obj{zero}))$ चे मूल्य (किंवा मध्यस्थ संकेतपद्धतीत,
$\Obj{two} \times (\Obj{three} + \Obj{zero})$) पुढीलप्रमाणे आहे:

\begin{multline*}
\Value{\times(\Obj{two}, +(\Obj{three},\Obj{zero}))}{M} \\
\begin{aligned}
& =\Assign{\times}{M}(\Value{\Obj{two}}{M}, \Value{+(\Obj{three}, \Obj{zero})}{M})\\
& = \Assign{\times}{M}(\Value{\Obj{two}}{M}, \Assign{+}{M}(\Value{\Obj{three}}{M},
\Value{\Obj{zero}}{M})) \\
& = \Assign{\times}{M}(\Assign{\Obj{two}}{M}, \Assign{+}{M}(\Assign{\Obj{three}}{M},
\Assign{\Obj{zero}}{M})) \\
& = \Assign{\times}{M}(2, \Assign{+}{M}(3, 0)) \\
& = \Assign{\times}{M}(2, 3) \\
& = 6
\end{aligned}
\end{multline*}
\end{ex}

\begin{prob}
अंकगणिताचे मानक प्रतिमान~$\Struct N$ आच्छादित आहे का? स्पष्टीकरण द्या.
\end{prob}












\section{संरचना मधील
  सूत्राची पूर्ति}\label{fol:syn:sat:sec}

\begin{explain}
एकीकडे पदे आणि सूत्रे यांसारख्या अभिव्यक्ती आणि दुसरीकडे
संरचना यांना जोडणाऱ्या मूलभूत संकल्पना म्हणजे पदाचे
\emph{मूल्य} आणि सूत्राची \emph{पूर्ति}. अनौपचारिकपणे,
पदाचे मूल्य म्हणजे संरचना मधील घटक असतो---पद
नुसते स्थिरचिन्ह असेल, तर त्या स्थिरचिन्हाला संरचना ने नेमून
दिलेली वस्तू हे त्याचे मूल्य असते; आणि ते फलने वापरून
बनवलेले असेल, तर पदातील स्थिरचिन्हांची मूल्ये आणि फलनचिन्हांना
नेमून दिलेली फलने यांपासून त्याचे मूल्य काढले जाते. विधेयांना
दिलेल्या अर्थनिर्धारणामुळे संरचनेच्या प्रांतात सूत्र सत्य
होत असेल, तर सूत्राची त्या संरचनांमध्ये \emph{पूर्ति
होते}. पूर्तिची ही संकल्पना विगमनाधारित रीतीने विनिर्दिष्ट केली जाते:
आण्विक सूत्रांची पूर्ति केव्हा होते हे संरचनेचे
विनिर्देशन थेट सांगते; आणि एखाद्या जटिल सूत्राची पूर्ति केव्हा
होते हे आपण तिचा मुख्य संयोजक किंवा संख्यापक आणि तिच्या तात्काळ
उपसूत्रांची पूर्ति होते की नाही यांवरून परिभाषित करतो.

इथे संख्यापकांची स्थिती थोडी अवघड आहे, कारण संख्यकीकृत सूत्र
च्या तात्काळ उपसूत्रामध्ये एक मुक्त चल असते आणि
संरचना ही चलांची मूल्ये विनिर्दिष्ट करत नाहीत.
ही अडचण हाताळण्यासाठी आपण \emph{चर-मूल्यांकने} ही आणतो आणि पूर्तिची
व्याख्या केवळ संरचनेच्या सापेक्ष न देता, संरचना आणि
चल चे मूल्यांकन या दोहोंच्या सापेक्ष देतो.
\end{explain}

\begin{defn}[चर-मूल्यांकन]
संरचना~$\Struct{M}$ हिच्यासाठीचे \emph{चर-मूल्यांकन}~$s$
म्हणजे प्रत्येक चल चे~$\Domain M$ मधील घटक कडे
प्रतिचित्रण करणारे फलन; म्हणजेच, $s\colon \Var \to \Domain M$.
\end{defn}

\begin{explain}
संरचना प्रत्येक स्थिरांक ला मूल्य नेमून देते, तर
चर-मूल्यांकन प्रत्येक चराला मूल्य नेमून देते. पण त्यांच्यापासून बनलेली
पदेही प्रांत मधील घटक ना नावे देण्यासाठी वापरायची आहेत.
त्यासाठी आपण पदांचे मूल्य विगमनाधारित रीतीने परिभाषित करतो.
स्थिरांक आणि चरांसाठी मूल्य अनुक्रमे संरचना किंवा
चर-मूल्यांकन विनिर्दिष्ट करते तेच असते; अधिक जटिल पदांसाठी ते
संरचना ने फलने ना नेमून दिलेली फलने वापरून पुनरावर्ती
रीतीने काढले जाते.
\end{explain}

\begin{defn}[पदांचे मूल्य]
$t$ हे भाषा~$\Lang L$ हिचे पद, $\Struct M$ ही~$\Lang L$ साठीची
संरचना आणि $s$ हे~$\Struct M$ साठीचे चल चे
मूल्यांकन असेल, तर \emph{मूल्य}~$\Value{t}{M}[s]$ पुढीलप्रमाणे
परिभाषित केले जाते:
\begin{enumerate}
\item \indcase{t}{c}{$\Value{\indfrm}{M}[s] = \Assign{\indcomplex}{M}$.}
\item \indcase{t}{x}{$\Value{\indfrm}{M}[s] = s(\indcomplex)$.}
\item \indcase{t}{\Atom{f}{t_1, \ldots, t_n}}{
\[\Value{\indfrm}{M}[s] = \Assign{f}{M}(\Value{t_1}{M}[s], \ldots,
\Value{t_n}{M}[s]).
\]}
\end{enumerate}
\end{defn}

\begin{defn}[$x$-पर्याय]
$s$ हे संरचना~$\Struct M$ साठीचे चल चे मूल्यांकन
असेल, तर~$s$ पासून फार तर~$x$ ला नेमून दिलेल्या मूल्यातच भिन्न असणारे
$\Struct M$ साठीचे कोणतेही चल चे मूल्यांकन~$s'$ याला~$s$ चा
\emph{$x$-पर्याय} म्हणतात. $s'$ हा~$s$ चा $x$-पर्याय असेल, तर आपण
$\varAssign{s'}{s}{x}$ असे लिहितो.
\end{defn}

\begin{explain}
मूल्यांकन~$s$ च्या $x$-पर्यायाने~$x$ ला काही वेगळेच नेमून दिले
\emph{पाहिजे} असे नाही, हे लक्षात घ्या. वस्तुतः, प्रत्येक मूल्यांकन
स्वतःचाच $x$-पर्याय मानला जातो.
\end{explain}

\begin{defn}
  $s$ हे संरचना~$\Struct M$ साठीचे चल चे मूल्यांकन
  आणि $m \in \Domain{M}$ असेल, तर मूल्यांकन~$\Subst{s}{m}{x}$ हे पुढील
  रीतीने परिभाषित केलेले चर-मूल्यांकन आहे:
  \[\Subst{s}{m}{x}(y) = \begin{cases}
    m & \text{जर } y \ident x\\
    s(y) & \text{अन्यथा}.
  \end{cases}\]
\end{defn}

दुसऱ्या शब्दांत, $\Subst{s}{m}{x}$ हा~$s$ चा असा विशिष्ट $x$-पर्याय
आहे, जो प्रांतातील घटक~$m$ हे~$x$ ला नेमून देतो आणि~$x$ शिवाय
इतर चले ना~$s$ ने नेमून दिलेल्याच गोष्टी नेमून देतो.

\begin{defn}[पूर्ति]
\label{fol:syn:sat:defn:satisfaction}
संरचना~$\Struct M$ मध्ये चल चे मूल्यांकन~$s$ याच्या
सापेक्ष सूत्र~$\varphi$ ची पूर्ति, संकेतांत $\Sat{M}{\varphi}[s]$, पुढीलप्रमाणे
पुनरावर्ती रीतीने परिभाषित केली जाते. (``$\Sat{M}{\varphi}[s]$ नाही'' या
अर्थाने आपण $\Sat/{M}{\varphi}[s]$ असे लिहितो.)
\begin{enumerate}
\item
  \indcase{\varphi}{\lfalse}{$\Sat/{M}{\indfrm}[s]$.}

\item
  \indcase{\varphi}{\ltrue}{$\Sat{M}{\indfrm}[s]$.}

\item \indcase{\varphi}{\Atom{R}{t_1, \dots, t_n}}{$\Sat{M}{\indfrm}[s]$
  तेव्हा आणि केवळ तेव्हाच, जेव्हा $\langle \Value{t_1}{M}[s], \dots,
  \Value{t_n}{M}[s] \rangle \in \Assign{R}{M}$.}

\item \indcase{\varphi}{\eq[t_1][t_2]}{$\Sat{M}{\indfrm}[s]$ तेव्हा आणि
  केवळ तेव्हाच, जेव्हा $\Value{t_1}{M}[s] = \Value{t_2}{M}[s]$.}

\item
  \indcase{\varphi}{\lnot \psi}{$\Sat{M}{\indfrm}[s]$ तेव्हा आणि केवळ तेव्हाच,
    जेव्हा $\Sat/{M}{\psi}[s]$.}

\item
  \indcase{\varphi}{(\psi \land \chi)}{$\Sat{M}{\indfrm}[s]$ तेव्हा आणि केवळ
    तेव्हाच, जेव्हा $\Sat{M}{\psi}[s]$ आणि $\Sat{M}{\chi}[s]$.}

\item
  \indcase{\varphi}{(\psi \lor \chi)}{$\Sat{M}{\indfrm}[s]$ तेव्हा आणि केवळ
    तेव्हाच, जेव्हा $\Sat{M}{\psi}[s]$ किंवा $\Sat{M}{\chi}[s]$ (किंवा
    दोन्ही).}

\item
  \indcase{\varphi}{(\psi \lif \chi)}{$\Sat{M}{\indfrm}[s]$ तेव्हा आणि केवळ
    तेव्हाच, जेव्हा $\Sat/{M}{\psi}[s]$ किंवा $\Sat{M}{\chi}[s]$ (किंवा
    दोन्ही).}

\item
  \indcase{\varphi}{(\psi \liff \chi)}{$\Sat{M}{\indfrm}[s]$ तेव्हा आणि केवळ
    तेव्हाच, जेव्हा एकतर $\Sat{M}{\psi}[s]$ आणि $\Sat{M}{\chi}[s]$ दोन्ही,
    किंवा $\Sat{M}{\psi}[s]$ व $\Sat{M}{\chi}[s]$ यांपैकी एकही नाही.}

\item
  \indcase{\varphi}{\lforall[x][\psi]}{$\Sat{M}{\indfrm}[s]$ तेव्हा आणि केवळ
    तेव्हाच, जेव्हा प्रत्येक घटक~$m \in \Domain M$ साठी
    $\Sat{M}{\psi}[\Subst{s}{m}{x}]$.}

\item
  \indcase{\varphi}{\lexists[x][\psi]}{$\Sat{M}{\indfrm}[s]$ तेव्हा आणि केवळ
  तेव्हाच, जेव्हा किमान एका घटक~$m \in \Domain M$ साठी
  $\Sat{M}{\psi}[\Subst{s}{m}{x}]$.}
\end{enumerate}
\end{defn}

\begin{explain}
शेवटच्या दोन उपवाक्यांत
चर-मूल्यांकने महत्त्वाची आहेत. ``प्रत्येक $m \in
\Domain{M}$ साठी $\Sat{M}{\psi(m)}$'' असे म्हणून आपण
$\lforall[x][\psi(x)]$ ची पूर्ति परिभाषित करू शकत नाही.
 ``किमान एका $m \in \Domain{M}$ साठी $\Sat{M}{\psi(m)}$''
असे म्हणून आपण $\lexists[x][\psi(x)]$ ची पूर्ति परिभाषित करू शकत नाही.
कारण $m \in \Domain M$ असेल, तर ते भाषेचे चिन्ह नसते आणि म्हणून
$\psi(m)$ हे सूत्र नसते (म्हणजे $\Subst{\psi}{m}{x}$ अपरिभाषित
असते).~$\Struct{M}$ मधील प्रत्येक घटक चे नाव देणारी
स्थिरांक किंवा पदे
आपल्याजवळ उपलब्ध आहेत असेही गृहीत धरता येत नाही, कारण संरचना
च्या व्याख्येत तशी कोणतीही अट नाही. प्रमाण भाषेतील स्थिरांकाचा
संच गणनीय अनंत आहे; म्हणून $\Domain{M}$ हा गणनीय नसेल,
तर प्रत्येक वस्तूला नाव देण्यासाठी पुरेशी स्थिरांक सुद्धा नसतात.

ही अडचण सोडवण्यासाठी आपण चल ची मूल्यांकने आणतो; त्यांमुळे
चरे प्रांतातील घटक शी थेट जोडता येतात. मग, उदाहरणार्थ,
$\lexists[x][\psi(x)]$ ची~$\Struct M$
मध्ये पूर्ति तेव्हा आणि केवळ तेव्हाच होते, जेव्हा
किमान एका $m \in \Domain{M}$ साठी $\Sat{M}{\psi(m)}$, असे न म्हणता आपण म्हणतो:
तिची~$\Struct M$ मध्ये~$s$ च्या \emph{सापेक्ष} पूर्ति तेव्हा आणि केवळ
तेव्हाच होते, जेव्हा $\psi(x)$ ची~$\Subst{s}{m}{x}$ च्या सापेक्ष
किमान एका $m \in \Domain M$ साठी पूर्ति होते.

\end{explain}

\begin{ex}
$\Lang{L} = \{a, b, f, R\}$ असू द्या, जिथे $a$ आणि $b$ ही
स्थिरांक, $f$ हे दोन-स्थानी फलन आणि $R$ हे दोन-स्थानी
विधेय आहे. पुढीलप्रमाणे परिभाषित केलेली संरचना~$\Struct{M}$
विचारात घ्या:
\begin{enumerate}
\item $\Domain M = \{1, 2, 3, 4\}$
\item $\Assign{a}{M} = 1$
\item $\Assign{b}{M} = 2$
\item $x+y \le 3$ असेल तर $\Assign{f}{M}(x, y) = x+y$, अन्यथा $= 3$.
\item $\Assign{R}{M} = \{\tuple{1, 1}, \tuple{1, 2}, \tuple{2, 3}, \tuple{2, 4}\}$
\end{enumerate}
प्रत्येक चल ला $1 \in \Domain{M}$ नेमून देणारे फलन
$s(x) = 1$ हे~$\Struct{M}$ साठीचे चर-मूल्यांकन आहे.

मग
\begin{align*}
\Value{f(a,b)}{M}[s] & = \Assign{f}{M}(\Value{a}{M}[s], \Value{b}{M}[s]).
\intertext{$a$ आणि $b$ ही स्थिरांक असल्यामुळे $\Value{a}{M}[s]
  = \Assign{a}{M} = 1$ आणि $\Value{b}{M}[s] = \Assign{b}{M} = 2$. म्हणून}
\Value{f(a,b)}{M}[s] & = \Assign{f}{M}(1, 2) = 1+2 = 3.
\intertext{$f(f(a,b),a)$ चे मूल्य काढण्यासाठी पुढील अभिव्यक्ती विचारात घ्यावी लागते:}
\Value{f(f(a,b),a)}{M}[s] & = \Assign{f}{M}(\Value{f(a, b)}{M}[s],
\Value{a}{M}[s]) = \Assign{f}{M}(3, 1) = 3,
\intertext{कारण $3+1 > 3$. $s(x) = 1$ आणि $\Value{x}{M}[s] =
  s(x)$ असल्यामुळे आपल्याकडे पुढील समीकरणही आहे:}
\Value{f(f(a,b),x)}{M}[s] & = \Assign{f}{M}(\Value{f(a, b)}{M}[s],
\Value{x}{M}[s]) = \Assign{f}{M}(3, 1) = 3,
\end{align*}

आण्विक सूत्र~$R(t_1, t_2)$ च्या फलसाधकांच्या मूल्यांची क्रमित
जोडी, म्हणजे $\tuple{\Value{t_1}{M}[s], \Value{t_2}{M}[s]}$, ही
$\Assign{R}{M}$ ची घटक असेल, तर त्या सूत्राची पूर्ति होते.
म्हणून, उदाहरणार्थ, $\tuple{\Value{b}{M}, \Value{f(a,b)}{M}} =
\tuple{2, 3} \in \Assign{R}{M}$ असल्यामुळे
$\Sat{M}{R(b,f(a,b))}[s]$; पण $\tuple{1, 3} \notin \Assign{R}{M}$
असल्यामुळे $\Sat/{M}{R(x, f(a,b))}[s]$.


अनाण्विक सूत्र~$\varphi$ ची पूर्ति होते का हे ठरवण्यासाठी त्याच्या मुख्य
संयोजकाला लागू होणारे विगमनाधारित व्याख्येतील उपवाक्य वापरावे. उदाहरणार्थ,
$R(a, a) \lif (R(b, x) \lor R(x, b))$ मधील मुख्य संयोजक~$\lif$ आहे आणि
\begin{align*}
  & \Sat{M}{R(a, a) \lif (R(b, x) \lor R(x, b))}[s] \text{ तेव्हा आणि केवळ तेव्हाच, जेव्हा }\\
  & \qquad
  \Sat/{M}{R(a,a)}[s] \text{ किंवा } \Sat{M}{R(b, x) \lor R(x, b)}[s]
  \intertext{$\Sat{M}{R(a,a)}[s]$ असल्यामुळे (कारण $\tuple{1,1} \in
    \Assign{R}{M}$), उत्तर अजून ठरवता येत नाही; प्रथम
    $\Sat{M}{R(b, x) \lor R(x, b)}[s]$ आहे का ते शोधावे लागते:}
  & \Sat{M}{R(b, x) \lor R(x, b)}[s] \text{ तेव्हा आणि केवळ तेव्हाच, जेव्हा }\\
  & \qquad \Sat{M}{R(b,x)}[s] \text{ किंवा } \Sat{M}{R(x, b)}[s]
  \intertext{आणि हे खरे आहे, कारण $\Sat{M}{R(x, b)}[s]$
    (कारण $\tuple{1,2} \in \Assign{R}{M}$).}
\end{align*}

आठवा, $s$ चा $x$-पर्याय म्हणजे~$s$ पासून फार तर~$x$ ला नेमून
दिलेल्या मूल्यातच भिन्न असणारे चर-मूल्यांकन. $\Domain{M}$ च्या प्रत्येक
घटक साठी~$s$ चा एक $x$-पर्याय आहे:
\begin{align*}
  s_1 & = \Subst{s}{1}{x}, &
  s_2 & = \Subst{s}{2}{x},\\
  s_3 & = \Subst{s}{3}{x}, &
  s_4 & = \Subst{s}{4}{x}.
\end{align*}
म्हणून, उदाहरणार्थ, $s_2(x) = 2$ आणि~$x$ शिवाय इतर सर्व चर~$y$ साठी
$s_2(y) = s(y) = 1$. $\Domain{M} = \{1, 2, 3, 4\}$ असल्यामुळे हेच
$\Struct{M}$ या रचनेसाठीचे~$s$ चे सर्व $x$-पर्याय आहेत. विशेषतः,
$s_1 = s$ हे लक्षात घ्या ($s$ हा नेहमीच स्वतःचा $x$-पर्याय असतो).

अस्तित्ववाची संख्यापक असलेल्या
  सूत्र~$\lexists[x][\varphi(x)]$ ची पूर्ति होते का हे ठरवण्यासाठी,
  किमान एका $m \in \Domain M$ साठी
  $\Sat{M}{\varphi(x)}[\Subst{s}{m}{x}]$ आहे का हे ठरवावे लागते. म्हणून,
  \[  \Sat{M}{\lexists[x][(R(b,x) \lor R(x,b))]}[s],
  \]
  कारण $\Sat{M}{R(b,x) \lor R(x, b)}[\Subst{s}{1}{x}]$
  ($\Subst{s}{3}{x}$ हेसुद्धा चालले असते). पण,
  \[  \Sat/{M}{\lexists[x][(R(b,x) \land R(x,b))]}[s]
  \]
  कारण कोणताही $m \in \Domain{M}$ निवडला, तरी
  $\Sat/{M}{R(b,x) \land R(x,b)}[\Subst{s}{m}{x}]$.

वैश्विक संख्यापक असलेल्या
  सूत्र~$\lforall[x][\varphi(x)]$ ची पूर्ति होते का हे ठरवण्यासाठी,
  प्रत्येक $m \in \Domain M$ साठी
  $\Sat{M}{\varphi(x)}[\Subst{s}{m}{x}]$ आहे का हे ठरवावे लागते. म्हणून,
  \[  \Sat{M}{\lforall[x][(R(x,a) \lif R(a,x))]}[s],
  \]
  कारण प्रत्येक $m \in \Domain M$ साठी
  $\Sat{M}{R(x,a) \lif R(a,x)}[\Subst{s}{m}{x}]$. $m = 1$ साठी
  $\Sat{M}{R(a,x)}[\Subst{s}{1}{x}]$ असल्यामुळे उत्तरवर्ती सत्य आहे;
  $m = 2$, $3$ आणि~$4$ साठी $\Sat/{M}{R(x,a)}[\Subst{s}{m}{x}]$
  असल्यामुळे पूर्ववर्ती असत्य आहे. पण,
  \[  \Sat/{M}{\lforall[x][(R(a,x) \lif R(x,a))]}[s]
  \]
  कारण $\Sat/{M}{R(a,x) \lif R(x,a)}[\Subst{s}{2}{x}]$ (कारण
  $\Sat{M}{R(a, x)}[\Subst{s}{2}{x}]$ आणि $\Sat/{M}{R(x,
  a)}[\Subst{s}{2}{x}]$).





अधिक जटिल उदाहरणासाठी पुढील सूत्र विचारात घ्या:
\[\lforall[x][(R(a,x) \lif \lexists[y][R(x,y)])].
\]
$\Sat/{M}{R(a,x)}[\Subst{s}{3}{x}]$ आणि
$\Sat/{M}{R(a,x)}[\Subst{s}{4}{x}]$ असल्यामुळे, सशर्त सूत्राच्या
उत्तरवर्तीची काळजी घ्यावी लागणारी रंजक प्रकरणे फक्त $m = 1$ आणि
$m = 2$ ही आहेत. $\Sat{M}{\lexists[y][R(x,y)]}[\Subst{s}{1}{x}]$ आहे
का? किमान एक $n \in \Domain M$ असा असेल की
$\Sat{M}{R(x,y)}[\Subst{\Subst{s}{1}{x}}{n}{y}]$, तर ते खरे आहे.
वस्तुतः, $n = 1$ घेतल्यास $\Subst{\Subst{s}{1}{x}}{n}{y} =
\Subst{s}{1}{y} = s$. $s(x) = 1$, $s(y) = 1$ आणि
$\tuple{1,1} \in \Assign{R}{M}$ असल्यामुळे उत्तर होकारार्थी आहे.


$\Sat{M}{\lexists[y][R(x,y)]}[\Subst{s}{2}{x}]$ आहे का हे ठरवण्यासाठी
आपल्याला चल ची मूल्यांकने
$\Subst{\Subst{s}{2}{x}}{n}{y}$ तपासावी लागतात. इथे $n = 1$ साठी हे
मूल्यांकन $s_2 = \Subst{s}{2}{x}$ आहे आणि ते $R(x,y)$ ची पूर्ति करत
नाही ($s_2(x) = 2$, $s_2(y) = 1$ आणि
$\tuple{2,1}\notin \Assign{R}{M}$). पण
$\Subst{\Subst{s}{2}{x}}{3}{y} = \Subst{s_2}{3}{y}$ विचारात घ्या.
$\tuple{2,3} \in \Assign{R}{M}$ असल्यामुळे
$\Sat{M}{R(x,y)}[\Subst{s_2}{3}{y}]$ आणि म्हणून
$\Sat{M}{\lexists[y][R(x,y)]}[s_2]$.

म्हणून, प्रत्येक $m \in \Domain M$ साठी एकतर
$\Sat/{M}{R(a,x)}[\Subst{s}{m}{x}]$ ($m = 3$, $4$ असेल तर) किंवा
$\Sat{M}{\lexists[y][R(x,y)]}[\Subst{s}{m}{x}]$ ($m = 1$, $2$ असेल तर),
आणि म्हणून
\[\Sat{M}{\lforall[x][(R(a,x) \lif \lexists[y][R(x,y)])]}[s].
\]
दुसरीकडे,
\[\Sat/{M}{\lexists[x][(R(a,x) \land \lforall[y][R(x,y)])]}[s].
\]
फक्त $m = 1$ आणि $m = 2$ यांच्यासाठी
$\Sat{M}{R(a,x)}[\Subst{s}{m}{x}]$. पण~$m$ च्या या दोन्ही मूल्यांपैकी
प्रत्येकासाठी एक $n \in \Domain M$ आहे---पहिल्या मूल्यासाठी $n = 4$
आणि दुसऱ्यासाठी $n = 1$---ज्यामुळे
$\Sat/{M}{R(x,y)}[\Subst{\Subst{s}{m}{x}}{n}{y}]$ आणि म्हणून $m = 1$
व $m = 2$ साठी
$\Sat/{M}{\lforall[y][R(x,y)]}[\Subst{s}{m}{x}]$. एकंदरीत, असा कोणताही
$m \in \Domain M$ नाही की
$\Sat{M}{R(a,x) \land \lforall[y][R(x,y)]}[\Subst{s}{m}{x}]$.


\end{ex}









\begin{prob}
एक स्थिरांक, एक एक-स्थानी फलन आणि एक दोन-स्थानी
विधेय असलेली $\Lang L = \{c, f, A\}$ असू द्या; आणि
संरचना~$\Struct{M}$ पुढीलप्रमाणे दिलेली असू द्या:
\begin{enumerate}
\item $\Domain M = \{1, 2, 3\}$
\item $\Assign{c}{M} = 3$
\item $\Assign{f}{M}(1) = 2, \Assign{f}{M}(2) = 3, \Assign{f}{M}(3) = 2$
\item $\Assign{A}{M} = \{\tuple{1, 2}, \tuple{2, 3}, \tuple{3, 3}\}$
\end{enumerate}
(a) सर्व चले~$v$ साठी $s(v) = 1$ असू द्या. पुढील विधान
खरे आहे की नाही हे शोधा:
\[\Sat{M}{\lexists[x][(A(f(z), c) \lif \lforall[y][(A(y, x) \lor A(f(y),
      x))])]}[s]
\]
तुमचे उत्तर का योग्य आहे ते स्पष्ट करा.

(b) ज्यात या सूत्राची पूर्ति होत नाही अशी वेगळी रचना आणि
चल चे मूल्यांकन द्या.
\end{prob}












\section{चर-मूल्यांकने}\label{fol:syn:ass:sec}

\begin{explain}
चल चे मूल्यांकन~$s$ हे \emph{प्रत्येक} चराला मूल्य देते---आणि
चरे अनंत असतात. अर्थात, हे आवश्यक नाही. आपण चल च्या
मूल्यांकनांनी सर्व चले ना मूल्ये नेमावीत अशी अट घालतो, कारण
त्यामुळे गोष्टी खूप सोप्या होतात. पद~$t$ चे मूल्य आणि सूत्र~$\varphi$
ची संरचनांमध्ये~$s$ च्या सापेक्ष पूर्ति होते की नाही, हे फक्त
$t$ मध्ये येणाऱ्या चले ना आणि~$\varphi$ च्या मुक्त चले ना
$s$ ने केलेल्या नेमणुकांवर अवलंबून असते. पुढील दोन विधानांचा आशय हाच
आहे. “अवलंबून असणे” ही कल्पना नेमकी करण्यासाठी आपण दोन गोष्टी दाखवतो:
$t$ मधील सर्व चरांवर एकमत असणारी कोणतीही दोन चर-मूल्यांकने तेच मूल्य
देतात; आणि दोन चर-मूल्यांकने~$\varphi$ च्या सर्व मुक्त चरांवर एकमत असतील, तर
$\varphi$ ची एकाच्या सापेक्ष पूर्ति होते तेव्हा आणि केवळ तेव्हाच तिची
दुसऱ्याच्या सापेक्ष पूर्ति होते.
\end{explain}

\begin{prop}\label{fol:syn:ass:prop:valindep}
पद~$t$ मधील चले ही $x_1$, \dots,~$x_n$ यांपैकी असतील आणि
$i = 1$, \dots,~$n$ साठी $s_1(x_i) = s_2(x_i)$ असेल, तर
$\Value{t}{M}[s_1] = \Value{t}{M}[s_2]$.
\end{prop}

\begin{proof}
$t$ च्या जटिलतेवरील विगमनाने सिद्ध करू. आधार बाबतीत~$t$ हे
स्थिरांक किंवा $x_1$, \dots,~$x_n$ यांपैकी एखादे चर असू शकते.
$t = c$ असेल, तर
$\Value{t}{M}[s_1] = \Assign{c}{M} = \Value{t}{M}[s_2]$. $t = x_i$
असेल, तर विधानाच्या गृहीतकानुसार $s_1(x_i) = s_2(x_i)$; म्हणून
$\Value{t}{M}[s_1] = s_1(x_i) = s_2(x_i) = \Value{t}{M}[s_2]$.

विगमन पायरीसाठी $t = \Atom{f}{t_1, \dots, t_k}$ असून दावा
$t_1$, \dots, $t_k$ साठी खरा आहे असे समजू. मग
\begin{align*}
  \Value{t}{M}[s_1] & = \Value{\Atom{f}{t_1, \dots, t_k}}{M}[s_1] \\
  & = \Assign{f}{M}(\Value{t_1}{M}[s_1], \dots, \Value{t_k}{M}[s_1]).
\intertext{$j = 1$, \dots,~$k$ साठी~$t_j$ मधील चले ही
    $x_1$, \dots,~$x_n$ यांपैकी आहेत. विगमन गृहीतकानुसार
    $\Value{t_j}{M}[s_1] = \Value{t_j}{M}[s_2]$. म्हणून,}
\Value{t}{M}[s_1] & = \Value{\Atom{f}{t_1, \dots, t_k}}{M}[s_1] \\
  & = \Assign{f}{M}(\Value{t_1}{M}[s_1], \dots, \Value{t_k}{M}[s_1]) \\
  & = \Assign{f}{M}(\Value{t_1}{M}[s_2], \dots, \Value{t_k}{M}[s_2]) \\
  & = \Value{\Atom{f}{t_1, \dots, t_k}}{M}[s_2] = \Value{t}{M}[s_2].
\end{align*}
\end{proof}

\begin{prop}\label{fol:syn:ass:prop:satindep}
$\varphi$ मधील मुक्त चले ही $x_1$, \dots,~$x_n$ यांपैकी असतील आणि
$i = 1$, \dots,~$n$ साठी $s_1(x_i) = s_2(x_i)$ असेल, तर
$\Sat{M}{\varphi}[s_1]$ तेव्हा आणि केवळ तेव्हाच, जेव्हा
$\Sat{M}{\varphi}[s_2]$.
\end{prop}

\begin{proof}
$\varphi$ च्या जटिलतेवर विगमन वापरू. आधार बाबतीत~$\varphi$ हे आण्विक असते आणि
$\varphi$ पुढीलपैकी असू शकते:
$\ltrue$,
$\lfalse$,
$k$-स्थानी विधेय~$R$ आणि पदे $t_1$, \dots,~$t_k$ यांसाठी
$\Atom{R}{t_1, \dots, t_k}$; किंवा पदे~$t_1$ व~$t_2$ यांसाठी
$\eq[t_1][t_2]$. शेवटच्या दोन बाबतींत आपण द्विशर्त ची फक्त
पुढची दिशा दाखवतो, कारण उलट दिशेची सिद्धता सममित आहे.

\begin{enumerate}
\item
  \indcase{\varphi}{\ltrue}{$\Sat{M}{\indfrm}[s_1]$ आणि
    $\Sat{M}{\indfrm}[s_2]$ दोन्ही.}

\item
  \indcase{\varphi}{\lfalse}{$\Sat/{M}{\varphi}[s_1]$ आणि
    $\Sat/{M}{\varphi}[s_2]$ दोन्ही.}

\item
  \indcase{\varphi}{\Atom{R}{t_1, \ldots, t_k}}{
    $\Sat{M}{\indfrm}[s_1]$ असे समजू. मग
    \[    \langle \Value{t_1}{M}[s_1], \ldots, \Value{t_k}{M}[s_1] \rangle
    \in \Assign{R}{M}.
    \]
    $i = 1$, \dots,~$k$ साठी \ref{fol:syn:ass:prop:valindep} नुसार
    $\Value{t_i}{M}[s_1] = \Value{t_i}{M}[s_2]$. म्हणून आपल्याकडे
    $\langle \Value{t_1}{M}[s_2], \ldots, \Value{t_k}{M}[s_2] \rangle
    \in \Assign{R}{M}$ हेसुद्धा आहे; आणि म्हणून
    $\Sat{M}{\indfrm}[s_2]$.
    }
\item
  \indcase{\varphi}{\eq[t_1][t_2]}{$\Sat{M}{\indfrm}[s_1]$ असे समजू.
    मग $\Value{t_1}{M}[s_1] = \Value{t_2}{M}[s_1]$. म्हणून,
    \begin{align*}
      \Value{t_1}{M}[s_2] & = \Value{t_1}{M}[s_1]
      & \text{(\ref{fol:syn:ass:prop:valindep} नुसार)} \\
      & = \Value{t_2}{M}[s_1]
      & \text{(कारण $\Sat{M}{\eq[t_1][t_2]}[s_1]$)}\\
      &= \Value{t_2}{M}[s_2]
      & \text{(\ref{fol:syn:ass:prop:valindep} नुसार),}
    \end{align*}
    म्हणून $\Sat{M}{\eq[t_1][t_2]}[s_2]$.}
\end{enumerate}

आता~$\varphi$ पेक्षा कमी जटिल असलेल्या सर्व सूत्रे~$\psi$ साठी
$\Sat{M}{\psi}[s_1]$ तेव्हा आणि केवळ तेव्हाच, जेव्हा
$\Sat{M}{\psi}[s_2]$, असे समजू. विगमन पायरी~$\varphi$ च्या मुख्य कारकाने
ठरणाऱ्या बाबींनुसार पुढे जाते. प्रत्येक बाबतीत आपण द्विशर्त ची
फक्त पुढची दिशा दाखवतो; उलट दिशेची सिद्धता सममित आहे. संख्यापकांच्या
बाबी सोडून इतर सर्व बाबतींत आपण~$\varphi$ च्या उप-सूत्रे~$\psi$ ना
विगमन गृहीतक लावतो. $\psi$ ची मुक्त चरे ही~$\varphi$ च्या मुक्त चरांपैकी
असतात. त्यामुळे $s_1$ व~$s_2$ ही~$\varphi$ च्या मुक्त चरांवर एकमत असतील,
तर ती~$\psi$ च्या मुक्त चरांवरही एकमत असतात आणि विगमन गृहीतक~$\psi$ ला
लागू होते.

\begin{enumerate}
\item

    \indcase{\varphi}{\lnot \psi}{$\Sat{M}{\indfrm}[s_1]$ असेल, तर
      $\Sat/{M}{\psi}[s_1]$; म्हणून विगमन गृहीतकानुसार
      $\Sat/{M}{\psi}[s_2]$ आणि त्यामुळे $\Sat{M}{\indfrm}[s_2]$.}

\item

    \indcase{\varphi}{\psi \land \chi}{$\Sat{M}{\indfrm}[s_1]$ असेल, तर
      $\Sat{M}{\psi}[s_1]$ आणि $\Sat{M}{\chi}[s_1]$; म्हणून विगमन
      गृहीतकानुसार $\Sat{M}{\psi}[s_2]$ आणि $\Sat{M}{\chi}[s_2]$.
      त्यामुळे $\Sat{M}{\indfrm}[s_2]$.}

\item

    \indcase{\varphi}{\psi \lor \chi}{$\Sat{M}{\indfrm}[s_1]$ असेल, तर
      $\Sat{M}{\psi}[s_1]$ किंवा $\Sat{M}{\chi}[s_1]$. विगमन गृहीतकानुसार
      $\Sat{M}{\psi}[s_2]$ किंवा $\Sat{M}{\chi}[s_2]$; म्हणून
      $\Sat{M}{\indfrm}[s_2]$.}

\item

    \indcase{\varphi}{\psi \lif \chi}{$\Sat{M}{\indfrm}[s_1]$ असेल, तर
    $\Sat/{M}{\psi}[s_1]$ किंवा $\Sat{M}{\chi}[s_1]$. विगमन गृहीतकानुसार
    $\Sat/{M}{\psi}[s_2]$ किंवा $\Sat{M}{\chi}[s_2]$; म्हणून
    $\Sat{M}{\indfrm}[s_2]$.}

\item

    \indcase{\varphi}{\psi \liff \chi}{$\Sat{M}{\indfrm}[s_1]$ असेल, तर एकतर
      $\Sat{M}{\psi}[s_1]$ आणि $\Sat{M}{\chi}[s_1]$, किंवा
      $\Sat/{M}{\psi}[s_1]$ आणि $\Sat/{M}{\chi}[s_1]$. विगमन गृहीतकानुसार
      एकतर $\Sat{M}{\psi}[s_2]$ आणि $\Sat{M}{\chi}[s_2]$, किंवा
      $\Sat/{M}{\psi}[s_2]$ आणि $\Sat/{M}{\chi}[s_2]$. दोन्हींपैकी
      कोणतीही बाब असली, तरी $\Sat{M}{\indfrm}[s_2]$.}

\item

    \indcase{\varphi}{\lexists[x][\psi]}{$\Sat{M}{\indfrm}[s_1]$ असेल, तर असा
      एक $m \in \Domain{M}$ असतो की $\Sat{M}{\psi}[\Subst{s_1}{m}{x}]$.
      $s_1' = \Subst{s_1}{m}{x}$ आणि $s_2' =
      \Subst{s_2}{m}{x}$ असू द्या. $\psi$ ची मुक्त चरे
      $x_1$, \dots, $x_n$ आणि~$x$ यांपैकी आहेत. $s_1'$ आणि~$s_2'$ हे
      अनुक्रमे $s_1$ व~$s_2$ यांचे $x$-पर्याय असल्यामुळे आणि गृहीतकानुसार
      $s_1(x_i) = s_2(x_i)$ असल्यामुळे $s_1'(x_i) = s_2'(x_i)$.
      $s_1'$ आणि~$s_2'$ आपण ज्या रीतीने परिभाषित केली आहे, तिच्यामुळे
      $s_1'(x) = s_2'(x) = m$. मग~$\psi$ आणि $s_1'$, $s_2'$ यांना विगमन
      गृहीतक लागू होते; म्हणून $\Sat{M}{\psi}[s_2']$. $s_2' =
      \Subst{s_2}{m}{x}$ असल्यामुळे असा एक $m \in \Domain{M}$ आहे की
      $\Sat{M}{\psi}[\Subst{s_2}{m}{x}]$; म्हणून
      $\Sat{M}{\indfrm}[s_2]$.}

\item

    \indcase{\varphi}{\lforall[x][\psi]}{$\Sat{M}{\indfrm}[s_1]$ असेल, तर
      प्रत्येक $m \in \Domain{M}$ साठी
      $\Sat{M}{\psi}[\Subst{s_1}{m}{x}]$. प्रत्येक
      $m \in \Domain{M}$ साठी $\Sat{M}{\psi}[\Subst{s_2}{m}{x}]$
      हेसुद्धा दाखवायचे आहे. म्हणून मनमाना $m \in \Domain{M}$ घेऊन
      $s_1' = \Subst{s_1}{m}{x}$ आणि $s_2' =
      \Subst{s_2}{m}{x}$ विचारात घ्या. आपल्याकडे
      $\Sat{M}{\psi}[s_1']$. $\psi$ ची मुक्त चरे $x_1$, \dots, $x_n$
      आणि~$x$ यांपैकी आहेत. $s_1'$ आणि~$s_2'$ हे अनुक्रमे $s_1$
      व~$s_2$ यांचे $x$-पर्याय असल्यामुळे आणि गृहीतकानुसार
      $s_1(x_i) = s_2(x_i)$ असल्यामुळे $s_1'(x_i) = s_2'(x_i)$.
      $s_1'$ आणि~$s_2'$ आपण ज्या रीतीने परिभाषित केली आहे, तिच्यामुळे
      $s_1'(x) = s_2'(x) = m$. मग~$\psi$ आणि~$s_1'$, $s_2'$ यांना विगमन
      गृहीतक लागू होते आणि आपल्याकडे $\Sat{M}{\psi}[s_2']$. हे प्रत्येक
      $m \in \Domain{M}$ ला लागू होते, म्हणजे सर्व~$m \in \Domain{M}$
      साठी $\Sat{M}{\psi}[\Subst{s_2}{m}{x}]$; म्हणून
      $\Sat{M}{\indfrm}[s_2]$.
      }
\end{enumerate}
विगमनाने आपल्याला पुढील निष्कर्ष मिळतो: $\varphi$ मधील मुक्त चले
ही $x_1$, \dots, $x_n$ यांपैकी असतील आणि $i = 1$, \dots,~$n$ साठी
$s_1(x_i)=s_2(x_i)$ असेल, तेव्हा $\Sat{M}{\varphi}[s_1]$ तेव्हा आणि केवळ
तेव्हाच, जेव्हा $\Sat{M}{\varphi}[s_2]$.
\end{proof}



\begin{explain}
वाक्ये मध्ये मुक्त चरे नसतात; म्हणून कोणतीही दोन चर-मूल्यांकने
कोणत्याही वाक्याच्या सर्व (शून्य) मुक्त चरांना त्याच गोष्टी नेमून देतात.
आत्ताच सिद्ध केलेल्या विधानाचा अर्थ मग असा होतो की वाक्याची
एखाद्या रचनेत चर-मूल्यांकनाच्या सापेक्ष पूर्ति होते की नाही, हे पूर्णपणे
त्या मूल्यांकनापासून स्वतंत्र असते. हे तथ्य आपण नोंदवू. पुढे येणारी
वाक्याची संरचना मधील पूर्तिची व्याख्या (चर-मूल्यांकनाचा
उल्लेख न करता) त्यामुळे समर्थित होते.
\end{explain}

\begin{cor}
\label{fol:syn:ass:cor:sat-sentence}
$\varphi$ हे वाक्य आणि $s$ हे चर-मूल्यांकन असेल, तर प्रत्येक
चर-मूल्यांकन~$s'$ साठी $\Sat{M}{\varphi}[s]$ तेव्हा आणि केवळ तेव्हाच,
जेव्हा $\Sat{M}{\varphi}[s']$.
\end{cor}

\begin{proof}
  $s'$ हे कोणतेही चर-मूल्यांकन असू द्या. $\varphi$ हे वाक्य
  असल्यामुळे त्यात मुक्त चरे नाहीत; म्हणून प्रत्येक चर-मूल्यांकन~$s'$
  हे~$\varphi$ च्या सर्व मुक्त चरांना~$s$ सारख्याच गोष्टी आपोआप नेमून देते.
  त्यामुळे \ref{fol:syn:ass:prop:satindep} ची अट पूर्ण होते आणि
  $\Sat{M}{\varphi}[s]$ तेव्हा आणि केवळ तेव्हाच, जेव्हा
  $\Sat{M}{\varphi}[s']$.
\end{proof}

\begin{defn}
\label{fol:syn:ass:defn:satisfaction}
$\varphi$ हे वाक्य असेल, तर संरचना~$\Struct M$ ही~$\varphi$ ची
\emph{पूर्ति करते}, $\Sat{M}{\varphi}$, असे आपण तेव्हा आणि केवळ तेव्हाच
म्हणतो, जेव्हा सर्व चर-मूल्यांकने~$s$ साठी $\Sat{M}{\varphi}[s]$.
\end{defn}

$\Sat{M}{\varphi}$ असेल, तर आपण सरळ \emph{$\varphi$ हे~$\Struct{M}$ मध्ये सत्य
आहे} असेही म्हणतो. पूर्तिची संकल्पना स्वतंत्र वाक्यांपासून
वाक्यांच्या संचांपर्यंत स्वाभाविकपणे विस्तारते.

\begin{defn}
\label{fol:syn:ass:defn:sat}
$\Gamma$ हा वाक्यांचा संच असेल, तर संरचना~$\Struct M$
ही~$\Gamma$ ची \emph{पूर्ति करते}, $\Sat{M}{\Gamma}$, असे आपण तेव्हा
आणि केवळ तेव्हाच म्हणतो, जेव्हा प्रत्येक $\varphi \in \Gamma$ साठी
$\Sat{M}{\varphi}$.

\end{defn}

\begin{prop}\label{fol:syn:ass:prop:sentence-sat-true}
  $\Struct{M}$ ही संरचना, $\varphi$ हे वाक्य आणि $s$ हे
  चर-मूल्यांकन असू द्या. $\Sat{M}{\varphi}$ तेव्हा आणि केवळ तेव्हाच,
  जेव्हा $\Sat{M}{\varphi}[s]$.
\end{prop}

\begin{proof}
सराव.
\end{proof}

\begin{prob}
\ref{fol:syn:ass:prop:sentence-sat-true} सिद्ध करा.
\end{prob}

\begin{prop}\label{fol:syn:ass:prop:sat-quant}
  $\varphi(x)$ मध्ये फक्त~$x$ मुक्त येतो आणि $\Struct M$ ही
  संरचना आहे असे समजा. मग:
  \begin{enumerate}
    \item $\Sat{M}{\lexists[x][\varphi(x)]}$ तेव्हा आणि केवळ
      तेव्हाच, जेव्हा किमान एका चर-मूल्यांकन~$s$ साठी
      $\Sat{M}{\varphi(x)}[s]$.
    \item $\Sat{M}{\lforall[x][\varphi(x)]}$ तेव्हा आणि केवळ
      तेव्हाच, जेव्हा सर्व चर-मूल्यांकने~$s$ साठी
      $\Sat{M}{\varphi(x)}[s]$.
\end{enumerate}
\end{prop}

\begin{proof}
सराव.
\end{proof}

\begin{prob}
\ref{fol:syn:ass:prop:sat-quant} सिद्ध करा.
\end{prob}

\begin{prob}
\DeclareRobustCommand{\VDash}{\mathrel{||}\joinrel\Relbar}
$\Lang L$ ही फलने नसलेली भाषा आहे असे समजा. संरचना~$\Struct M$,
स्थिरांक~$c$ आणि $a \in \Domain M$ दिले असता, $\Struct M[a/c]$
ही~$\Struct M$ सारखीच संरचना आहे, फक्त
$\Assign{c}{M[a/c]} = a$, अशी व्याख्या करा. वाक्ये~$\varphi$ साठी
$\Struct M \VDash \varphi$ ची पुढीलप्रमाणे व्याख्या करा:
\begin{enumerate}
\item
  \indcase{\varphi}{\lfalse}{$\Struct{M} \VDash \indfrm$ नाही.}

\item
  \indcase{\varphi}{\ltrue}{$\Struct{M} \VDash \indfrm$.}

\item \indcase{\varphi}{\Atom{R}{d_1, \dots, d_n}}{$\Struct M \VDash \indfrm$
  तेव्हा आणि केवळ तेव्हाच, जेव्हा
  $\langle \Assign{d_1}{M}, \dots, \Assign{d_n}{M} \rangle \in
  \Assign{R}{M}$.}

\item \indcase{\varphi}{\eq[d_1][d_2]}{$\Struct M \VDash \indfrm$ तेव्हा
  आणि केवळ तेव्हाच, जेव्हा
  $\Assign{d_1}{M} = \Assign{d_2}{M}$.}

\item
  \indcase{\varphi}{\lnot \psi}{$\Struct{M} \VDash \indfrm$ तेव्हा आणि केवळ
    तेव्हाच, जेव्हा $\Struct{M} \VDash {\psi}$ नाही.}

\item
  \indcase{\varphi}{(\psi \land \chi)}{$\Struct{M} \VDash
    \indfrm$ तेव्हा आणि केवळ तेव्हाच, जेव्हा
    $\Struct{M} \VDash\psi$ आणि $\Struct{M} \VDash \chi$.}

\item
  \indcase{\varphi}{(\psi \lor \chi)}{$\Struct{M} \VDash \indfrm$ तेव्हा आणि
    केवळ तेव्हाच, जेव्हा $\Struct{M} \VDash \psi$ किंवा
    $\Struct{M} \VDash \chi$ (किंवा दोन्ही).}

\item
  \indcase{\varphi}{(\psi \lif \chi)}{$\Struct{M} \VDash \indfrm$ तेव्हा आणि
    केवळ तेव्हाच, जेव्हा $\Struct {M} \VDash \psi$ नाही किंवा
    $\Struct M \VDash \chi$ (किंवा दोन्ही).}

\item
  \indcase{\varphi}{(\psi \liff \chi)}{$\Struct{M} \VDash {\indfrm}$ तेव्हा आणि
    केवळ तेव्हाच, जेव्हा एकतर $\Struct{M} \VDash {\psi}$ आणि
    $\Struct{M} \VDash {\chi}$ दोन्ही, किंवा $\Struct{M} \VDash {\psi}$ व
    $\Struct{M} \VDash {\chi}$ यांपैकी एकही नाही.}

\item
  \indcase{\varphi}{\lforall[x][\psi]}{$\Struct{M} \VDash {\indfrm}$ तेव्हा आणि
    केवळ तेव्हाच, जेव्हा प्रत्येक $a \in \Domain{M}$ साठी
    $\Struct{M[a/c]} \VDash \Subst{\psi}{c}{x}$, जर~$\psi$ मध्ये~$c$ येत
    नसेल.}

\item
  \indcase{\varphi}{\lexists[x][\psi]}{$\Struct{M} \VDash
  {\indfrm}$ तेव्हा आणि केवळ तेव्हाच, जेव्हा असा एक $a \in \Domain M$
  आहे की $\Struct{M[a/c]} \VDash \Subst{\psi}{c}{x}$, जर~$\psi$ मध्ये~$c$
  येत नसेल.}
\end{enumerate}
$x_1$, \dots, $x_n$ ही~$\varphi$ मधील सर्व मुक्त चले,
$c_1$, \dots, $c_n$ ही~$\varphi$ मध्ये नसलेली स्थिरचिन्हे,
$a_1$, \dots, $a_n \in \Domain M$ आणि $s(x_i) = a_i$ असू द्या.

$\Sat{M}{\varphi}[s]$ तेव्हा आणि केवळ तेव्हाच, जेव्हा
$\Struct M[a_1/c_1,\dots,a_n/c_n]
\VDash \Subst{\Subst{\varphi}{c_1}{x_1}\dots}{c_n}{x_n}$, हे दाखवा.

(या प्रश्नावरून असे दिसते की प्रथम-क्रम तर्कशास्त्राची चिन्हार्थमीमांसा
चर-मूल्यांकनांविना देणे शक्य आहे.)
\end{prob}

\begin{prob}
$f$ हे~$\varphi(x,y)$ मध्ये नसलेले फलनचिन्ह आहे असे समजा. अशी
संरचना~$\Struct{M}$ आहे की
$\Sat{M}{\lforall[x][\lexists[y][\varphi(x,y)]]}$, तेव्हा आणि केवळ तेव्हाच
अशी~$\Struct M'$ आहे की
$\Sat{M'}{\lforall[x][\varphi(x,f(x))]}$, हे दाखवा.

(हा प्रश्न स्कोलेमचे प्रमेय म्हणून परिचित असलेल्या प्रमेयाची एक विशेष
बाब आहे; $\lforall[x][\varphi(x,f(x))]$ याला
$\lforall[x][\lexists[y][\varphi(x,y)]]$ चे \emph{स्कोलेम सामान्य रूप}
म्हणतात.)
\end{prob}












\section{विस्तारात्मकता}\label{fol:syn:ext:sec}

\begin{explain}
विस्तारात्मकता, जिला कधीकधी संबद्धता म्हणतात, अनौपचारिकपणे पुढीलप्रमाणे
व्यक्त करता येते: संरचना~$\Struct M$ मध्ये चल च्या
मूल्यांकन~$s$ च्या सापेक्ष सूत्र~$\varphi$ ची पूर्ति होण्यावर परिणाम
करणारे घटक म्हणजे फक्त प्रांताचा आकार आणि~$\varphi$ मध्ये प्रत्यक्ष
येणाऱ्या भाषेच्या घटकांना~$\Struct M$ व~$s$ यांनी केलेल्या नेमणुका.

विस्तारात्मकतेचा एक तात्काळ परिणाम असा आहे: दोन संरचना
$\Struct M$ आणि~$\Struct M'$ यांचा प्रांत एकच असेल आणि वाक्य~$\varphi$ मध्ये
येणाऱ्या भाषेच्या सर्व घटकांबाबत त्या एकमत असतील, तर~$\varphi$ स्वतः सत्य
आहे की नाही याबाबतही $\Struct M$ आणि~$\Struct M'$ एकमत असले पाहिजेत.
\end{explain}

\begin{prop}[विस्तारात्मकता]
  \label{fol:syn:ext:prop:extensionality}
  $\varphi$ हे सूत्र, $\Struct M_1$ व~$\Struct M_2$ या
  $\Domain{M_1} = \Domain{M_2}$ असलेल्या संरचना, आणि~$s$ हे
  $\Domain{M_1} = \Domain{M_2}$ वरील चर-मूल्यांकन असू द्या. $\varphi$ मध्ये
  येणारे प्रत्येक स्थिरांक~$c$, संबंधचिन्ह~$R$ आणि फलन~$f$
  यांसाठी $\Assign{c}{M_1} = \Assign{c}{M_2}$,
  $\Assign{R}{M_1}=\Assign{R}{M_2}$ आणि
  $\Assign{f}{M_1} = \Assign{f}{M_2}$ असेल, तर
  $\Sat{M_1}{\varphi}[s]$ तेव्हा आणि केवळ तेव्हाच, जेव्हा
  $\Sat{M_2}{\varphi}[s]$.
\end{prop}

\begin{proof}
  प्रथम (पद~$t$ वरील विगमनाने) प्रत्येक पदासाठी
  $\Value{t}{M_1}[s] = \Value{t}{M_2}[s]$ हे सिद्ध करा. मग नुकताच
  सिद्ध केलेला दावा विगमनाच्या आधार बाबतीत ($\varphi$ आण्विक असताना) वापरून,
  $\varphi$ वरील विगमनाने विधान सिद्ध करा.
\end{proof}

\begin{prob}
\ref{fol:syn:ext:prop:extensionality} ची सिद्धता सविस्तर पूर्ण करा.
\end{prob}

\begin{cor}[वाक्ये साठी विस्तारात्मकता]
  \label{fol:syn:ext:cor:extensionality-sent}
  $\varphi$ हे वाक्य आणि $\Struct{M_1}$, $\Struct{M_2}$ या
  \ref{fol:syn:ext:prop:extensionality} मधीलप्रमाणे असू द्या. मग
  $\Sat{M_1}{\varphi}$ तेव्हा आणि केवळ तेव्हाच, जेव्हा $\Sat{M_2}{\varphi}$.
\end{cor}

\begin{proof}
\ref{fol:syn:ext:prop:extensionality} आणि \ref{fol:syn:ass:cor:sat-sentence} यांवरून
हे निष्पन्न होते.
\end{proof}

शिवाय, पदाचे मूल्य आणि संरचना ही सूत्राची पूर्ति करते की
नाही, हे फक्त तिच्या उपपदांच्या मूल्यांवर अवलंबून असते.

\begin{prop}\label{fol:syn:ext:prop:ext-terms}
$\Struct M$ ही संरचना, $t$ व~$t'$ ही पदे आणि $s$ हे
चर-मूल्यांकन असू द्या. मग $\Value{\Subst{t}{t'}{x}}{M}[s] =
\Value{t}{M}[\Subst{s}{\Value{t'}{M}[s]}{x}]$.
\end{prop}

\begin{proof}
$t$ वरील विगमनाने सिद्ध करू.
\begin{enumerate}
\item $t$ हे स्थिरचिन्ह, म्हणजे $t\ident c$, असेल, तर
  $\Subst{t}{t'}{x} = c$ आणि $\Value{c}{M}[s] = \Assign{c}{M} =
  \Value{c}{M}[\Subst{s}{\Value{t'}{M}[s]}{x}]$.

\item $t$ हे~$x$ शिवाय दुसरे चर, म्हणजे $t \ident y$, असेल, तर
  $\Subst{t}{t'}{x} = y$ आणि
  $\varAssign{s}{\Subst{s}{\Value{t'}{M}[s]}{x}}{x}$ असल्यामुळे
  $\Value{y}{M}[s] =
  \Value{y}{M}[\Subst{s}{\Value{t'}{M}[s]}{x}]$.

\item $t \ident x$ असेल, तर $\Subst{t}{t'}{x} = t'$. पण
  $\Subst{s}{\Value{t'}{M}[s]}{x}$ च्या व्याख्येनुसार
  $\Value{x}{M}[\Subst{s}{\Value{t'}{M}[s]}{x}] = \Value{t'}{M}[s]$.

\item $t \ident \Atom{f}{t_1,\dots,t_n}$ असेल, तर:
\begin{multline*}
  \Value{\Subst{t}{t'}{x}}{M}[s] \\
  \begin{aligned}[b]
& = \Value{\Atom{f}{\Subst{t_1}{t'}{x}, \dots, \Subst{t_n}{t'}{x}}}{M}[s]\\
& \qquad    \text{$\Subst{t}{t'}{x}$ च्या व्याख्येनुसार}\\
& = \Assign{f}{M}(\Value{\Subst{t_1}{t'}{x}}{M}[s], \dots,
    \Value{\Subst{t_n}{t'}{x}}{M}[s])\\
    & \qquad  \text{$\Value{\Atom{f}{\dots}}{M}[s]$ च्या व्याख्येनुसार}\\
& = \Assign{f}{M}(\Value{t_1}{M}[\Subst{s}{\Value{t'}{M}[s]}{x}], \dots,
   \Value{t_n}{M}[\Subst{s}{\Value{t'}{M}[s]}{x}])\\
& \qquad    \text{विगमन गृहीतकानुसार}\\
& = \Value{t}{M}[\Subst{s}{\Value{t'}{M}[s]}{x}]
    \text{$\Value{\Atom{f}{\dots}}{M}[\Subst{s}{\Value{t'}{M}[s]}{x}]$ च्या व्याख्येनुसार}
  \end{aligned}
\end{multline*}

\end{enumerate}
\end{proof}

\begin{prop}\label{fol:syn:ext:prop:ext-formulas} $\Struct M$ ही
संरचना, $\varphi$ हे सूत्र, $t'$ हे~A मध्ये~x साठी
आदेशनासाठी मुक्त असलेले पद आणि $s$ हे चर-मूल्यांकन असू द्या. मग
$\Sat{M}{\Subst{\varphi}{t'}{x}}[s]$ तेव्हा आणि केवळ तेव्हाच, जेव्हा
$\Sat{M}{\varphi}[\Subst{s}{\Value{t'}{M}[s]}{x}]$.

\end{prop}

\begin{proof}
सराव.
\end{proof}

\begin{prob}
\ref{fol:syn:ext:prop:ext-formulas} सिद्ध करा.
\end{prob}

\begin{explain}
  \ref{fol:syn:ext:prop:ext-terms} आणि \ref{fol:syn:ext:prop:ext-formulas} यांचा
  आशय पुढीलप्रमाणे आहे. आपल्याजवळ पद~$t$ किंवा सूत्र~$\varphi$ आणि
  दुसरे पद~$t'$ असून, $\Subst{t}{t'}{x}$ चे मूल्य किंवा
  $\Subst{\varphi}{t'}{x}$ ची संरचना~$\Struct M$ मध्ये
  चल च्या मूल्यांकन~$s$ च्या सापेक्ष पूर्ति होते की नाही हे
  जाणून घ्यायचे आहे असे समजा. मग आपण प्रथम आदेशन करून नंतर
  $\Struct{M}$ आणि~$s$ यांच्या सापेक्ष मूल्य किंवा पूर्ति विचारात घेऊ
  शकतो; किंवा प्रथम~$\Struct{M}$ मध्ये~$s$ च्या सापेक्ष~$t'$ चे
  मूल्य~$m = \Value{t'}{M}[s]$ ठरवून, चल चे मूल्यांकन
  बदलून~$\Subst{s}{m}{x}$ करू शकतो, आणि नंतर~$\Struct{M}$ व
  $\Subst{s}{m}{x}$ यांच्या सापेक्ष~$t$ चे मूल्य किंवा
  $\Sat{M}{\varphi}[\Subst{s}{m}{x}]$ आहे की नाही हे विचारात घेऊ शकतो.
  आपण यांपैकी कोणतीही पद्धत वापरली, तरी उत्तर एकच असेल, याची हमी
  \ref{fol:syn:ext:prop:ext-terms} आणि \ref{fol:syn:ext:prop:ext-formulas}
  देतात.
\end{explain}











\section{चिन्हार्थविषयक संकल्पना}\label{fol:syn:sem:sec}

\begin{explain}
प्रथम-क्रम भाषांसाठीच्या संरचनेची व्याख्या मिळाल्यावर वाक्यांचे
काही मूलभूत चिन्हार्थविषयक गुणधर्म आणि वाक्यांमधील संबंध परिभाषित करता
येतात. त्यांतील सर्वांत सोपी संकल्पना म्हणजे वाक्याची \emph{वैधता}.
प्रत्येक संरचनांमध्ये एखाद्या वाक्याची पूर्ति होत असेल, तर ते
वाक्य वैध असते. त्यातील अतार्किक चिन्हांचा अर्थ कसा लावला आहे याचा
विचार न करता ज्या वाक्यांची पूर्ति होते, तीच वैध वाक्ये होत. म्हणून वैध
वाक्यांना \emph{तार्किक सत्ये} असेही म्हणतात---ती कोणत्याही
संरचनांमध्ये सत्य, म्हणजे पूर्त, असतात. त्यामुळे त्यांची सत्यता
फक्त त्यांत येणारी तार्किक चिन्हे आणि त्यांची विन्यासात्मक संरचना
यांवर अवलंबून असते; अतार्किक चिन्हे किंवा त्यांचे अर्थनिर्धारण यांवर नाही.
\end{explain}

\begin{defn}[वैधता]
वाक्य~$\varphi$ हे \emph{वैध}, $\Entails \varphi$, तेव्हा आणि केवळ तेव्हाच असते,
जेव्हा प्रत्येक संरचना~$\Struct M$ साठी $\Sat{M}{\varphi}$.
\end{defn}

\begin{defn}[तार्किक निष्पन्नता]
वाक्यांच्या संच~$\Gamma$ मधून वाक्य~$\varphi$
\emph{तार्किकरीत्या निष्पन्न होते}, $\Gamma \Entails \varphi$, तेव्हा आणि
केवळ तेव्हाच, जेव्हा $\Sat{M}{\Gamma}$ असलेल्या प्रत्येक
संरचना~$\Struct M$ साठी $\Sat{M}{\varphi}$.
\end{defn}


\begin{defn}[पूर्ततायोग्यता]
एखाद्या संरचना~$\Struct M$ साठी $\Sat{M}{\Gamma}$ असेल, तर
वाक्यांचा संच~$\Gamma$ \emph{पूर्ततायोग्य} असतो. $\Gamma$
पूर्ततायोग्य नसेल, तर त्याला \emph{अपूर्ततायोग्य} म्हणतात.
\end{defn}

\begin{prop}
वाक्य~$\varphi$ वैध असते तेव्हा आणि केवळ तेव्हाच, जेव्हा
वाक्यांच्या प्रत्येक संच~$\Gamma$ साठी $\Gamma \Entails \varphi$.
\end{prop}

\begin{proof}
अग्र दिशेसाठी, $\varphi$ वैध आणि $\Gamma$ हा वाक्यांचा संच असू द्या.
$\Sat{M}{\Gamma}$ अशी संरचना~$\Struct M$ असू द्या. $\varphi$ वैध
असल्यामुळे $\Sat{M}{\varphi}$; म्हणून $\Gamma \Entails \varphi$.

उलट दिशेच्या व्यत्यासासाठी, $\varphi$ अवैध असू द्या; म्हणजे
$\Sat/{M}{\varphi}$ अशी संरचना~$\Struct M$ आहे. $\Gamma =
\{ \ltrue \}$ असताना $\ltrue$ वैध असल्यामुळे $\Sat{M}{\Gamma}$.
त्यामुळे अशी संरचना~$\Struct M$ आहे की $\Sat{M}{\Gamma}$ पण
$\Sat/{M}{\varphi}$; म्हणून $\Gamma$ मधून~$\varphi$ तार्किकरीत्या निष्पन्न होत
नाही.
\end{proof}

\begin{prop}
\label{fol:syn:sem:prop:entails-unsat}
$\Gamma \Entails \varphi$ तेव्हा आणि केवळ तेव्हाच, जेव्हा
$\Gamma \cup \{\lnot \varphi\}$ अपूर्ततायोग्य असतो.
\end{prop}

\begin{proof}
अग्र दिशेसाठी, $\Gamma \Entails \varphi$ असे समजू आणि, याउलट, अशी
संरचना~$\Struct M$ आहे की $\Sat{M}{\Gamma
  \cup \{ \lnot \varphi \}}$, असे समजू. $\Sat{M}{\Gamma}$ आणि
$\Gamma \Entails \varphi$ असल्यामुळे $\Sat{M}{\varphi}$. तसेच
$\Sat{M}{\Gamma\cup \{ \lnot \varphi \}}$ असल्यामुळे
$\Sat{M}{\lnot \varphi}$; म्हणून $\Sat{M}{\varphi}$ आणि $\Sat/{M}{\varphi}$ दोन्ही
मिळतात, हा विरोधाभास आहे. त्यामुळे अशी कोणतीही
संरचना~$\Struct M$ असू शकत नाही; म्हणून
$\Gamma \cup \{ \lnot \varphi \}$ अपूर्ततायोग्य आहे.

उलट दिशेसाठी, $\Gamma \cup \{ \lnot \varphi \}$ अपूर्ततायोग्य आहे असे
समजू. म्हणून प्रत्येक संरचना~$\Struct M$ साठी एकतर
$\Sat/{M}{\Gamma}$ किंवा $\Sat{M}{\varphi}$. त्यामुळे
$\Sat{M}{\Gamma}$ असलेल्या प्रत्येक संरचना~$\Struct M$ साठी
$\Sat{M}{\varphi}$; म्हणून $\Gamma \Entails \varphi$.
\end{proof}

\begin{prob}
\begin{enumerate}
\item $\Gamma \Entails \bot$ तेव्हा आणि केवळ तेव्हाच, जेव्हा $\Gamma$
  अपूर्ततायोग्य आहे, हे दाखवा.
\item $\Gamma \cup \{\varphi\} \Entails \bot$ तेव्हा आणि केवळ तेव्हाच,
  जेव्हा $\Gamma \Entails \lnot \varphi$, हे दाखवा.
\item $\varphi$ किंवा $\Gamma$ मध्ये~$c$ येत नाही असे समजा.
  $\Gamma \Entails \lforall[x][\varphi]$ तेव्हा आणि केवळ तेव्हाच, जेव्हा
  $\Gamma \Entails \Subst{\varphi}{c}{x}$, हे दाखवा.
\end{enumerate}
\end{prob}

\begin{prop}
$\Gamma \subseteq \Gamma'$ आणि $\Gamma \Entails \varphi$ असेल, तर
$\Gamma' \Entails \varphi$.
\end{prop}

\begin{proof}
$\Gamma \subseteq \Gamma'$ आणि $\Gamma \Entails \varphi$ असे समजा.
$\Sat{M}{\Gamma'}$ अशी रचना~$\Struct M$ असू द्या; मग
$\Sat{M}{\Gamma}$ आणि $\Gamma \Entails \varphi$ असल्यामुळे
$\Sat{M}{\varphi}$. म्हणून जेव्हा जेव्हा $\Sat{M}{\Gamma'}$, तेव्हा
$\Sat{M}{\varphi}$; त्यामुळे $\Gamma' \Entails \varphi$.
\end{proof}


\begin{thm}[चिन्हार्थविषयक निगमन प्रमेय]
\label{fol:syn:sem:thm:sem-deduction}
$\Gamma \cup \{\varphi\} \Entails \psi$ तेव्हा आणि केवळ तेव्हाच, जेव्हा
$\Gamma \Entails \varphi \lif \psi$.
\end{thm}

\begin{proof}
अग्र दिशेसाठी, $\Gamma \cup \{ \varphi \} \Entails \psi$ असू द्या आणि
$\Sat{M}{\Gamma}$ अशी संरचना~$\Struct M$ असू द्या.
$\Sat{M}{\varphi}$ असेल, तर $\Sat{M}{\Gamma \cup \{ \varphi \} }$; आणि
$\Gamma \cup \{ \varphi \}$ मधून~$\psi$ तार्किकरीत्या निष्पन्न होत असल्यामुळे
$\Sat{M}{\psi}$. म्हणून $\Sat{M}{\varphi \lif \psi}$; त्यामुळे
$\Gamma \Entails \varphi \lif \psi$.

उलट दिशेसाठी, $\Gamma \Entails \varphi \lif \psi$ असू द्या आणि
$\Sat{M}{\Gamma \cup \{ \varphi \}}$ अशी संरचना~$\Struct M$ असू
द्या. मग $\Sat{M}{\Gamma}$; म्हणून $\Sat{M}{\varphi \lif \psi}$ आणि
$\Sat{M}{\varphi}$ असल्यामुळे $\Sat{M}{\psi}$. त्यामुळे जेव्हा जेव्हा
$\Sat{M}{\Gamma \cup \{ \varphi \} }$, तेव्हा $\Sat{M}{\psi}$; म्हणून
$\Gamma \cup \{ \varphi \} \Entails \psi$.
\end{proof}

\begin{prop}\label{fol:syn:sem:prop:quant-terms}
  $\Struct{M}$ ही संरचना, $\varphi(x)$ हे एक मुक्त चर~$x$ असलेले
  सूत्र आणि $t$ हे बंद पद असू द्या. मग:
  \begin{enumerate}
    \item $\varphi(t) \Entails \lexists[x][\varphi(x)]$
    \item $\lforall[x][\varphi(x)] \Entails \varphi(t)$
  \end{enumerate}
\end{prop}

\begin{proof}
  \begin{enumerate}
    \item
      $\Sat{M}{\varphi(t)}$ असे समजा. $s(x) =
        \Value{t}{M}$ असे चर-मूल्यांकन~$s$ असू द्या. $\varphi(t)$ हे
        वाक्य असल्यामुळे $\Sat{M}{\varphi(t)}[s]$.
        \ref{fol:syn:ext:prop:ext-formulas} नुसार
        $\Sat{M}{\varphi(x)}[s]$. \ref{fol:syn:ass:prop:sat-quant} नुसार
        $\Sat{M}{\lexists[x][\varphi(x)]}$.
    \item
      $\Sat{M}{\lforall[x][\varphi(x)]}$ असे समजा.
        $s(x) = \Value{t}{M}$ असे चर-मूल्यांकन~$s$ असू द्या.
        \ref{fol:syn:ass:prop:sat-quant} नुसार $\Sat{M}{\varphi(x)}[s]$.
        \ref{fol:syn:ext:prop:ext-formulas} नुसार
        $\Sat{M}{\varphi(t)}[s]$. $\varphi(t)$ हे वाक्य असल्यामुळे
        \ref{fol:syn:ass:prop:sentence-sat-true} नुसार
        $\Sat{M}{\varphi(t)}$.
  \end{enumerate}
\end{proof}





